% GENERATED FILE. Edit canonical lessons, then run npm run generate:books.
% canonical-source-hash: fnv1a64:2f84a284fb57bf42
% canonical-lessons: FA-C368-tahlil, FA-C368-tahlil-kardan, FA-C368-bardasht, FA-C368-natije-giri, FA-C368-eshare-kardan

\chapter{Looking closely}
\label{ch:fa-looking-closely}
\hlchaptermodality{368}

\begin{chapteropening}
\textbf{By the end of this chapter:} \emph{I can analyse something, give my impression of it, draw a conclusion and mention a point.}

Complete the last lesson of chapter 368: I can analyse something, give my impression of it, draw a conclusion and mention a point.
\end{chapteropening}

\section[\texorpdfstring{\fa{تحلیل} (tahlil)}{تحلیل (tahlil)}]{\fa{تحلیل} (tahlil) — an analysis}

\label{lesson:FA-C368-tahlil}



\begin{quote}
Before the new one: say the Persian for nevertheless, then the Persian for although.
\end{quote}

\subsection*{You'll want to know: \fa{تحلیل}}
\textbf{\fa{تحلیل}} (\emph{tahlil}) — \textquotedblleft{}an analysis\textquotedblright{}.

\begin{quote}
\fa{تحلیل} \fa{او} \fa{از} \fa{بازی} \fa{خیلی} \fa{خوب} \fa{بود}.
\end{quote}

\begin{quote}
\emph{tahlil-e u az bâzi kheyli khub bud.}
\end{quote}

\begin{quote}
His analysis of the game was very good.
\end{quote}

Plural \emph{tahlilhâ}. \emph{Tahlilgar} — an analyst.

\begin{cousinweb}
\textbf{\fa{تحلیل}} is usually traced to Arabic \emph{taḥlīl} \textquotedblleft{}analysis\textquotedblright{}, a verbal noun usually explained as untying or breaking something into its parts.
\end{cousinweb}

\begin{grammarlens}[title={What you've built}]
One more word for telling, explaining and describing at B1.
\end{grammarlens}

\subsection*{Your turn}
\begin{itemize}
\raggedright
  \item \emph{Say it:} \emph{tahlil}
  \item \emph{Say it:} \emph{tahlil}, once more
  \item \emph{Say it:} a sentence of your own with \emph{tahlil}
  \item \emph{Recall:} say the Persian for nevertheless, then the Persian for although, then say \emph{tahlil} again
\end{itemize}

\begin{tcolorbox}[breakable,colback=teal!4,colframe=teal!35!black,title={Before you move on}]
What does \fa{تحلیل} mean? (\textquotedblleft{}An analysis\textquotedblright{}.) Say it once more.
\end{tcolorbox}
\section[\texorpdfstring{\fa{تحلیل} \fa{کردن} (tahlil kardan)}{تحلیل کردن (tahlil kardan)}]{\fa{تحلیل} \fa{کردن} (tahlil kardan) — to analyse}

\label{lesson:FA-C368-tahlil-kardan}



\begin{quote}
Before the new one: say the Persian for although, then the Persian for an analysis.
\end{quote}

\subsection*{You'll want to know: \fa{تحلیل} \fa{کردن}}
\textbf{\fa{تحلیل} \fa{کردن}} (\emph{tahlil kardan}) — \textquotedblleft{}to analyse\textquotedblright{}.

\begin{quote}
\fa{استاد} \fa{نتیجه} \fa{آزمایش} \fa{را} \fa{تحلیل} \fa{کرد}.
\end{quote}

\begin{quote}
\emph{ostâd natije-ye âzmâyesh râ tahlil kard.}
\end{quote}

\begin{quote}
The professor analysed the result of the experiment.
\end{quote}

\emph{Tahlil} + \emph{kardan} (stem \emph{kon-}). The object takes \emph{râ}.

\begin{cousinweb}
\textbf{\fa{تحلیل} \fa{کردن}} joins \emph{tahlil}, usually traced to Arabic \emph{taḥlīl} \textquotedblleft{}analysis\textquotedblright{}, to the inherited Persian verb \emph{kardan} \textquotedblleft{}to do\textquotedblright{}.
\end{cousinweb}

\begin{grammarlens}[title={What you've built}]
One more word for telling, explaining and describing at B1.
\end{grammarlens}

\subsection*{Your turn}
\begin{itemize}
\raggedright
  \item \emph{Say it:} \emph{tahlil kardan}
  \item \emph{Say it:} \emph{tahlil kardan}, once more
  \item \emph{Say it:} a sentence of your own with \emph{tahlil kardan}
  \item \emph{Recall:} say the Persian for although, then the Persian for an analysis, then say \emph{tahlil kardan} again
\end{itemize}

\begin{tcolorbox}[breakable,colback=teal!4,colframe=teal!35!black,title={Before you move on}]
What does \fa{تحلیل} \fa{کردن} mean? (\textquotedblleft{}To analyse\textquotedblright{}.) Say it once more.
\end{tcolorbox}
\section[\texorpdfstring{\fa{برداشت} (bardâsht)}{برداشت (bardâsht)}]{\fa{برداشت} (bardâsht) — an impression, an interpretation}

\label{lesson:FA-C368-bardasht}



\begin{quote}
Before the new one: say the Persian for an analysis, then the Persian for to analyse.
\end{quote}

\subsection*{You'll want to know: \fa{برداشت}}
\textbf{\fa{برداشت}} (\emph{bardâsht}) — \textquotedblleft{}an impression, an interpretation\textquotedblright{}.

\begin{quote}
\fa{برداشت} \fa{من} \fa{این} \fa{بود} \fa{که} \fa{او} \fa{ناراحت} \fa{است}.
\end{quote}

\begin{quote}
\emph{bardâsht-e man in bud ke u nârâhat ast.}
\end{quote}

\begin{quote}
My impression was that he was upset.
\end{quote}

The past stem of \emph{bardâshtan} \textquotedblleft{}to pick up\textquotedblright{}: what you take from something. It is also a harvest.

\begin{cousinweb}
\textbf{\fa{برداشت}} is native Persian, built from \emph{bar} \textquotedblleft{}up, on\textquotedblright{} + \emph{dâsht}, the past stem of \emph{dâshtan} \textquotedblleft{}to have, to hold\textquotedblright{}. It is what you pick up from a text or a meeting.
\end{cousinweb}

\begin{grammarlens}[title={What you've built}]
One more word for telling, explaining and describing at B1.
\end{grammarlens}

\subsection*{Your turn}
\begin{itemize}
\raggedright
  \item \emph{Say it:} \emph{bardâsht}
  \item \emph{Say it:} \emph{bardâsht}, once more
  \item \emph{Say it:} a sentence of your own with \emph{bardâsht}
  \item \emph{Recall:} say the Persian for an analysis, then the Persian for to analyse, then say \emph{bardâsht} again
\end{itemize}

\begin{tcolorbox}[breakable,colback=teal!4,colframe=teal!35!black,title={Before you move on}]
What does \fa{برداشت} mean? (\textquotedblleft{}An impression, an interpretation\textquotedblright{}.) Say it once more.
\end{tcolorbox}
\section[\texorpdfstring{\fa{نتیجه} \fa{گیری} (natije giri)}{نتیجه گیری (natije giri)}]{\fa{نتیجه} \fa{گیری} (natije giri) — a conclusion (drawn)}

\label{lesson:FA-C368-natije-giri}



\begin{quote}
Before the new one: say the Persian for to analyse, then the Persian for an impression.
\end{quote}

\subsection*{You'll want to know: \fa{نتیجه} \fa{گیری}}
\textbf{\fa{نتیجه} \fa{گیری}} (\emph{natije giri}) — \textquotedblleft{}a conclusion (drawn)\textquotedblright{}.

\begin{quote}
\fa{نتیجه} \fa{گیری} \fa{او} \fa{منطقی} \fa{بود}.
\end{quote}

\begin{quote}
\emph{natije giri-ye u manteqi bud.}
\end{quote}

\begin{quote}
His conclusion was logical.
\end{quote}

\emph{Natije} \textquotedblleft{}result\textquotedblright{} + \emph{giri} \textquotedblleft{}taking\textquotedblright{}, from \emph{gereftan}. Also the last part of an essay.

\begin{cousinweb}
\textbf{\fa{نتیجه} \fa{گیری}} joins \emph{natije}, from Arabic \emph{natīja} \textquotedblleft{}result\textquotedblright{} (EtymDB), to \emph{giri} \textquotedblleft{}taking\textquotedblright{}, built on \emph{gir-}, the present stem of the inherited Persian verb \emph{gereftan} \textquotedblleft{}to take\textquotedblright{}.
\end{cousinweb}

\begin{grammarlens}[title={What you've built}]
One more word for telling, explaining and describing at B1.
\end{grammarlens}

\subsection*{Your turn}
\begin{itemize}
\raggedright
  \item \emph{Say it:} \emph{natije giri}
  \item \emph{Say it:} \emph{natije giri}, once more
  \item \emph{Say it:} a sentence of your own with \emph{natije giri}
  \item \emph{Recall:} say the Persian for to analyse, then the Persian for an impression, then say \emph{natije giri} again
\end{itemize}

\begin{tcolorbox}[breakable,colback=teal!4,colframe=teal!35!black,title={Before you move on}]
What does \fa{نتیجه} \fa{گیری} mean? (\textquotedblleft{}A conclusion (drawn)\textquotedblright{}.) Say it once more.
\end{tcolorbox}
\section[\texorpdfstring{\fa{اشاره} \fa{کردن} (eshâre kardan)}{اشاره کردن (eshâre kardan)}]{\fa{اشاره} \fa{کردن} (eshâre kardan) — to mention, to refer to}

\label{lesson:FA-C368-eshare-kardan}



\begin{quote}
Before the new one: say the Persian for an impression, then the Persian for a conclusion.
\end{quote}

\subsection*{You'll want to know: \fa{اشاره} \fa{کردن}}
\textbf{\fa{اشاره} \fa{کردن}} (\emph{eshâre kardan}) — \textquotedblleft{}to mention, to refer to\textquotedblright{}.

\begin{quote}
\fa{او} \fa{در} \fa{نامه} \fa{به} \fa{این} \fa{مشکل} \fa{اشاره} \fa{کرد}.
\end{quote}

\begin{quote}
\emph{u dar nâme be in moshkel eshâre kard.}
\end{quote}

\begin{quote}
In the letter he mentioned this problem.
\end{quote}

Takes \emph{be}: \emph{be chizi eshâre kardan}. Literally to make a sign towards something.

\begin{cousinweb}
\textbf{\fa{اشاره}} was borrowed from Arabic \emph{ishāra} \textquotedblleft{}a sign, a gesture\textquotedblright{} (EtymDB); with the Persian verb \emph{kardan} it means to refer to something, in a gesture or in words.
\end{cousinweb}

\begin{grammarlens}[title={What you've built}]
One more word for telling, explaining and describing at B1.
\end{grammarlens}

\subsection*{Your turn}
\begin{itemize}
\raggedright
  \item \emph{Say it:} \emph{eshâre kardan}
  \item \emph{Say it:} \emph{eshâre kardan}, once more
  \item \emph{Say it:} a sentence of your own with \emph{eshâre kardan}
  \item \emph{Recall:} say the Persian for an impression, then the Persian for a conclusion, then say \emph{eshâre kardan} again
\end{itemize}

\begin{tcolorbox}[breakable,colback=teal!4,colframe=teal!35!black,title={Before you move on}]
What does \fa{اشاره} \fa{کردن} mean? (\textquotedblleft{}To mention, to refer to\textquotedblright{}.) Say it once more.
\end{tcolorbox}
